% Detailed derivation from the reviewed unit-pair and equality-wide records.
\section{Proofs of the joint cutoff estimates}
\label{pca:joint-details}

Write \([F]_p=R^{-1/p}\|F\|_{\ell^p}\) for counting norms on the
actual rows and \([F]_\infty=\|F\|_\infty\). Use the same convention
after fixing the second pivot, since \(A\asymp R\) and there are
only boundedly many second pivots. The physical sieve factors are
\[
 \begin{aligned}
 \Phi_6(A,B)&=A+B+A^{5/6}B^{1/3}+A^{1/3}B^{5/6},\\
 \Phi_3(A,B)&=A+B+2(AB)^{2/3},\\
 \Delta_6(x)&=\max(0,x-3/4,x/3-1/8,5x/6-1/2),\\
 \Delta_3(x)&=\max(0,x-3/4,2x/3-1/4).
 \end{aligned}
\]
The last two functions record the power of the corresponding
sieve divided by \(A\) at \(B=D^x\). We display bounds up to
\((DH)^\epsilon\) and the appropriate \(W\) seminorms.
The dyads, divisor multiplicities and harmonic norm sums have
subpower cost. Shortened lengths are replaced by one only in
the sieve factor, with a bounded nonempty subunit counted directly.

\subsection{Exact cutoff arrays and the equality moment}
\label{pca:joint-arrays}

For Theorem~\ref{pca:joint-full}, fix the ordered \(d=mn\).
On the support \(N\zeta_i>L_i\), the primal unused good
squarefree ideal \(c\) has \(Nc<2D/(EN\zeta)<7\). Thus the
partner primal unit part is the sum with \(c=1\). Its coefficient
is \(\mu(d)\mathbf1_{(d,u)=1}\) times the product of the two
Möbius and raw factors at \(\zeta_1,\zeta_2\), with the common
weight \(W(EN\zeta_1N\zeta_2/D)\).
It keeps \((d,\zeta)=1\), both original upper bounds \(N\zeta_i\le U\),
and \((\zeta_1,\zeta_2)=1\). Each nonzero summand has
\(N\zeta_i<7L_i\). Expand the common-prime test using
\(\sum_{w\mid(\zeta_1,\zeta_2)}\mu(w)\), with \(w\) good
squarefree and coprime to \(d\), and set \(\zeta_i=w b_i\).
The residual \(b_i\) avoid \(w,d\). Keep \(\mu(w)\psi_u(w)^2\)
and the row deletion. Logarithmic Fourier separation of \(W\)
gives two arrays \(B_s,B_l\) with lengths \(M_i=L_i/Nw\),
where \(s,l\) are the shorter and longer cutoffs. Their
coefficients retain the exact endpoints. In particular their
true masses multiply to \(U/(\sqrt E Nw)\).

For the equality vector, at fixed \(e_1,g,\ell\) and a dyad
\(P\le Nh<2P\), group the actual \(h=re_0\) coefficients
over the permissible \(e_0\)-allocations. The row factor of the
resulting sum \(F_P\) is \(\psi_u(h)^2\). The coefficient
is divisor bounded with the original middle and cutoff predicates.
The external inverse norm is \(D/(E E_1NgN\ell\sqrt R)\),
where \(E_1=Ne_1\) and the actual primitive factor \(Q_u^{-1/2}\)
is bounded using \(Q_u\asymp R\).
The Gauss root and fixed raw factors are retained.
Put \(T_0=\sqrt{D/(ER)}\), \(T_1=T_0/\sqrt{E_1}\).
On a dyad \(P\le Nh<2P\), the normalized weight from
the inverse square and the radial log separation is
\[
 P^{-2}\bigl(1+(T_1/P)^2\bigr)^{-4}\mathcal M_{12}(W).
\]
This keeps \(Q_u\) in the separated phase.

For a fixed integer \(j\ge2\), take the \(j\)-th power of the
formed \(h\)-array and group by the product of its \(j\) ideals.
The fixed factors on the second pivot and on \(S\) stay in the
coefficients. In the variable first pivot the \(h\)-character is
cubic. Write the grouped column as \(b_0q^3\), with \(b_0\)
cubefree, and freeze \(q\) with its raw multiplier
\(\psi_u(q)^6\). The squared mass on a fibre is bounded by
\(P^j/(Nq)^3\) up to divisor factors. The sum of the square roots
converges in \(q\), and the cubic comparison gives
\([F_P]_{2j}\ll \sqrt P\,(\Phi_3(A,P^j)/A)^{1/(2j)}\).
When \(T_1\ge1\) the radial dyad sum is bounded at \(P\asymp T_1\):
below it use the eighth power in the weight, and above it the
sieve factor grows at most as \(P^{1/2}\) against \(P^{-3/2}\).
When \(T_1<1\) use unit scale. For \(e\le1/4\), take
\(\tau=(1/4-e)/2\); the replacement of \(T_1\) by \(T_0\)
in the sieve costs only the monotone enlargement, and
\(E_1^{-1}T_1^{-3/2}=E_1^{-1/4}T_0^{-3/2}\).
In the subunit case \(E_1^{-1}\le T_0^{-2}\) when \(T_0\ge1\).
The equality bound after the \(g,\ell\) harmonic sums is
\begin{equation}\label{pca:joint-equality-moment}
 [\mathcal E_d]_{2j}\ll
 \sqrt{D/E}\,D^{-\tau/2+\Delta_3(j\tau)/(2j)}
       (DH)^\epsilon\mathcal M_{12}(W).
\end{equation}
The same dyad calculation with the elementary ideal count
gives the pointwise factor
\[
 \sqrt{D/E}\min(1,\sqrt{D/(ER)}).
\]

Set \(l_i=\log_D\max(1,L_i)\). The cutoff second moment has
excess \(\Delta_6(l_i)/2\) over its square-root mass.
For the fourth moment, square the actual cutoff sum.
The grouped product is sixth-free, with exponent two
allowed at a common prime, and has divisor-bounded
multiplicity. Its fourth-root cost is
\(D^{\Delta_6(2l_i)/4}\).
For \(e\le1/8\), pair the equality sixth moment with the
shorter third and longer second moments. Interpolating the
shorter between second and fourth with weight \(2/3\) on
the fourth yields
\[
 \frac{\Delta_6(l_l)}2+\frac{\Delta_6(l_s)}6
   +\frac{\Delta_6(2l_s)}6-\frac{\tau}2
   +\frac{\Delta_3(3\tau)}6
 =\frac1{48}-\frac{l_l-l_s}{12}.
\]
Here \(3/8\le l_s\le l_l\le1/2\), \(l_s+l_l=1-e\) at the level
of powers, and \(\Delta_3(3\tau)=0\).
For \(e\ge1/8\) the pointwise equality bound and the two
second moments have excess at most \(1/48\), because
\((l_s/3-1/8)_++(l_l/3-1/8)_+\le1/24\).
The normalized masses and \(w\)-sum give \(R/E\).
The original \(\mathcal E_d^{\rm unit}\) compensation has
normalized energy of order at most \(R/E\) by the grouped
elementary equality bound, and so does \(\mathcal E_d\).
Their inner product has the same upper bound. Summing over
\(d\) harmonically proves the \(37/48\) comparison.

\subsection{A forced radical in the physical sieve}
\label{pca:joint-radical}

For an actual good sixth-free ideal \(\mathfrak b\) of norm
of order \(B\), put
\(\mathfrak b^\sharp=\mathfrak b/\operatorname{rad}\mathfrak b\).
If every column satisfies \(N\operatorname{rad}(\mathfrak b^\sharp)
\gtrsim P\), the physical positive comparison can use
\begin{equation}\label{pca:joint-psi}
 \begin{aligned}
 \widetilde\Psi(A,B,P)={}&A+B+A^{5/6}B^{1/3}P^{-1/3}\\
              &+A^{1/6}BP^{-1/2}+A^{2/3}B^{1/2}P^{-1/2}.
 \end{aligned}
\end{equation}
The factor multiplies the original squared coefficient mass,
up to subpowers. Here is the support restriction inside
the pinned de Faveri argument. In its Lemma 4.4 remove the
outer common ideal of norm \(F\). The remaining scales obey
\(B'\asymp B/F\) and \(A/(CF)\le A'\le CA\) for a fixed constant
\(C\), and the radical threshold can drop at most by \(F\).
A prime lost from that radical divides the removed ideal.
The sharp part has norm at least the norm of its radical.

In the two positive norms of the proof of Lemma 5.1,
fix the sharp label \(s_i\) and the internal Möbius
divisor \(r\), with squarefree part \(r u_i\).
The original coefficient already requires
\(\operatorname{rad}(s_i)\mid r u_i\).
Set \(q_i=\operatorname{rad}(s_i)/(\operatorname{rad}(s_i),r)\),
so \(u_i=q_i v_i\). The \(v_i\) are squarefree and avoid \(r,q_i\).
The new label is in bijection with the original label at
these fixed indices. Keep the original denominator at
\(r q_i v_i\), the ray predicates, and the raw factor
\(\chi_{q_i}\) at the grouped row variable. Multiplicativity
gives the same zeros, including when a row prime divides
\(q_i\). This factor is a contraction in the formed norm.
The original squared coefficient sum is unchanged by the reindexing.

If \(P_i\) and \(C_i\) are the norm scales of \(s_i\) and
its radical, the row cap is still \(K_i\lesssim (B')^2/(A'P_i)\).
The column cap is now \(B'/(P_i C_i)\), since
\(N(rq_i)\ge N\operatorname{rad}(s_i)\).
Both \(P_i,C_i\) are at least of order \(\max(1,P/F)\).
Insert the four terms of the sixth-power comparison from
de Faveri's Theorem 1.1 and Lemma 6.1,
\(K+K^{1/6}L+K^{5/6}L^{1/3}+K^{1/3}L^{5/6}\).
The two factors in its positive Cauchy estimate carry
the prefactor \(A'\sqrt{P_1P_2}/B'\). Consequently
the four monomials before taking their geometric mean
are bounded by
\[
 \begin{aligned}
 B',\quad &(A')^{5/6}(B')^{1/3}P_i^{-1/6}C_i^{-1},\\
 &(A')^{1/6}B'P_i^{-1/6}C_i^{-1/3},\\
 &(A')^{2/3}(B')^{1/2}P_i^{-1/6}C_i^{-5/6}.
 \end{aligned}
\]
Put \(G=\max(1,P/F)\). The resulting factors satisfy
\[
 \begin{aligned}
 F^{-1/3}G^{-7/6}&\le P^{-1/3},\\
 F^{-1}G^{-1/2}&\le P^{-1/2},\\
 F^{-1/2}G^{-1}&\le P^{-1/2}.
 \end{aligned}
\]
These follow by considering \(F\le P\) and
\(F\ge P\). The source diagonal term gives \(A\).
This proves \eqref{pca:joint-psi} with its full masks.

For the pure joint norm apply Cauchy with the longer cutoff
as one factor. The other norm, of \(F_P\overline{B_s}\),
equals that of \(F_P B_s\) by pointwise modulus.
The exact signed pairing keeps its original conjugation.
The latter positive norm groups actual \(h^2 b_s\) injectively:
exponents \(0,1,2,3\) determine membership of \(h,b_s\)
at each prime. The radical of its sharp part is \(h\).
The length is of order \(P^2\max(1,M_s)\) and the squared
mass is bounded by \(P M_s\) up to subpowers. The square root
of \eqref{pca:joint-psi} grows at most linearly in \(P\)
at fixed \(M_s\); the inverse square and square-root mass
supply \(P^{-3/2}\). Thus the same radial maximum at
\(T_0\), with the \(E_1\)-factor as above, applies.
For \(e\le1/4\) the joint excess over the square-root mass
is \(\tfrac12(l_s/3+\tau/3-1/8)_+\); the \(B\) term is
below \(A\) in power, and the other two terms in
\eqref{pca:joint-psi} have no larger exponent.
With the longer second moment the total relative power is
\[
 \tfrac12(l_s/3+\tau/3-1/8)_+
 +\tfrac12(l_l/3-1/8)_+-\tau/2
     =\tfrac16(1/4+e/2-l_s)_+ .
\]
The equality follows from \(l_l\ge3/8\) and \(l_s+l_l=1-e\).
If \(e\le1/6\), then \(l_s\ge1/2-e\ge1/4+e/2\).
The middle compensation still costs \(R/E\), proving the
low-divisor and length claims of Theorem~\ref{pca:joint-full}.

\subsection{The selected Poisson product}
\label{pca:joint-selected-proof}

For Theorem~\ref{pca:joint-modes}, keep the actual selected
partner frequencies in the Poisson expression. The live dyad
gives \(N\zeta_i<14L_i\), and \(E_1Nh<14\) on the partner
wide support. Keep the ordered \(e_0\)-assignment, the
common-prime projector and all shortened residual deletions.
The actual partner frequency weight separates in logarithmic
variables on a frequency dyad, with absolute Fourier mass
bounded by \(\mathcal M_{12}(W)(1+T/R)^{-4}\).
The primitive factor remains \(Q_u^{-1/2}\) and the factor
from the re-expressed direct term remains in the array.

At fixed equality \(g,\ell,e_1\) and \(P\le Nh_0<2P\),
the raw equality phase times the conjugated partner dual
phase is \(\psi_u(h_0^2\mathfrak k)\). The row-independent
selector is included in the coefficient at fixed
\(h_0,\mathfrak k\) by summing the exact equality allocations.
It can give a nonseparable coefficient, with divisor-bounded
multiplicity and all the original signs. Write the grouped
product uniquely as \(\mathfrak s b_0 v^6\), with
\(\mathfrak s\) supported on \(S\), \(v\) good and \(b_0\)
good sixth-free. Factor \(h_0=h_{\rm in}h_{\rm out}\)
by support in \(v\). Then \(h_{\rm in}\mid v\) and
\(h_{\rm out}^2\mid b_0\), even when the actual frequency
meets \(h_0\).
At fixed \(\mathfrak s,v,h_{\rm in}\), count \(h_{\rm out}\)
at length \(P/Nh_{\rm in}\) and divide the physical residual
length by \((Nh_{\rm out})^2\).
The squared coefficient mass is bounded by
\[
 (DH)^\epsilon\,\frac{PT\,Nh_{\rm in}}{N\mathfrak s\,(Nv)^6}.
\]
The exact multipliers at \(\mathfrak s,v\) remain in the
positive comparison. The sums of the square-root weights
converge, since the sum over \(h_{\rm in}\mid v\) costs at
most \((Nv)^{1/2+\epsilon}\).
Also \(N\operatorname{rad}(b_0^\sharp)\ge P/Nv\).
Apply \eqref{pca:joint-psi} with this threshold, truncated
below by one, and with the true length at most of order
\(P^2T/(N\mathfrak s(Nv)^6)\).
Each inverse threshold exponent \(\beta\) is at most six
times its length exponent \(\alpha\). Thus the sieve may
use the nominal arguments \(P^2T,P\) after this shortening.
The radial sum at fixed \(T\) still has its maximum at the
equality scale, as the sieve square root grows at most
linearly in \(P\).

Put \(\Gamma_4=(\Delta_6(2l_s)+\Delta_6(2l_l))/4\).
The two fourth moments of the cutoffs give this excess,
with their true shortened masses.
Before the sieve factors, the normalized mass at fixed labels
and radial dyad, omitting the harmless fixed shortening and
the harmonic common-prime factor, is
\[
 \frac{\sqrt{L_1L_2}\sqrt T\,P^{-3/2}}{E E_1 NgN\ell}.
\]
The factor \(P^{-3/2}\) includes the equality inverse square
and its square-root mass. The outer normalization by \(D\)
cancels the equality prefactor \(D\); conversion of the
counting norm back to the row sum supplies \(R\) against
the two factors of size \(Q_u^{-1/2}\).
For \(e\le1/4\) and the indicated frequency ranges above
\(t=3/8\), the relative joint sieve excess is
\((t/3+\tau/3-1/8)/2\). Adding the physical frequency factor,
the cutoff moments and the equality radial gain gives
\(2t/3-7/16+\Gamma_4-\tau/3\).
For \(e\le1/8\), \(\Gamma_4=1/6-5e/12\).
For \(1/8\le e\le1/4\), convexity and the individual length
caps give \(\Gamma_4\le13/96-e/6\).
These are the first two branches of
\eqref{pca:joint-frequency-frontier}; smaller \(t\) follows
by monotonicity of the upper comparison.

For \(e\ge1/4\), use the unit-scale equality dyad and the
unshortened frequency sieve at \(T\). The total power above
\(R/E\) is
\[
 t/2-1/4-e/2+\Gamma+\Delta_6(t)/2 .
\]
For the third and fourth branches one may take
\(\Gamma=\Gamma_4\), bounded by \(13/96-e/6\)
through \(e=5/16\) and by \(1/12\) thereafter.
For \(e\ge27/56\) use the lesser of \(\Gamma_4\) and
\(\Gamma_2=l_s/2+\Delta_6(l_l)/2\); the latter
uses the shorter pointwise bound and longer second moment.
Under \(l_s+l_l\le29/56\) and \(l_l\le1/2\),
the maximum of this minimum is \(3/56\).
Indeed \(l_l\) may be increased to its permitted endpoint;
the two affine bounds cross at \(l_s=5/56,\ l_l=3/7\).
For \(l_s\le1/56\) the pointwise option is smaller than
\(3/56\), and for \(l_s\ge1/7\) convexity bounds the
fourth-moment option by its values at \(1/7,29/112\).
At \(t=3/4\) the resulting power is at most
\(3/16-e/2+3/56\le0\).
Above \(R\), the remaining growth has slope at most one
and the factor \((1+T/R)^{-4}\) gives a summable tail.
For equality-dependent selectors, sum the actual equality
label vectors only after using each label's own
selected middle equality subtraction. Each of these
subtractions has the elementary normalized energy \(R/E\).
The sum of the equality-label pointwise majorants is
\(O((DH)^\epsilon\mathcal M_{12}(W)\sqrt{D/E})\), by the
original harmonic \(g,\ell\) sums and the inverse-square radial
bound. Cauchy at each label gives the required total
compensation of order \(R/E\). The divisor sum is harmonic.

\subsection{Sparse frequency and equality support}
\label{pca:joint-sparse-proof}

For a frequency-only selector, rational content at least \(D^b\)
leaves \(O(TD^{-b})\) ideals on a dyad. Divide each frequency
by its rational principal factor and sum the inverse squares
of the possible contents. For \(j=1,2\), power the formed
frequency array \(j\) times before taking its second moment.
Grouping by product gives squared mass at most
\((DH)^\epsilon(TD^{-b})^j\) by the divisor bound.
At fixed \(S\)-part \(\mathfrak s\), a dyad \(Q\) of the
sixth-power base \(q\) has two bounds for the sum of square-root
masses. They are \(T^{j/2} Q^{1/2}D^{-jb/2}\)
by Cauchy over disjoint fibres, and \(T^{j/2}Q^{-2}\)
by the individual shortened length and ideal count.
The sieve length is at most of order \(\max(1,T^j/Q^6)\).
On writing \(x=\log_D Q/j\), the bound is convex up to \(b/5\)
and nonincreasing afterwards. Its maximum is bounded by the
larger value at \(0,b/5\). Thus the \(2j\)-norm of the
frequency array has excess
\[
 \lambda_j=
 \max\left\{\frac{\Delta_6(jt)}{2j}-\frac b2,\
   \frac{\Delta_6(j(t-6b/5)_+)}{2j}-\frac{2b}5\right\}.
\]
The two cutoff moments and \eqref{pca:joint-equality-moment}
give the following upper powers above \(R/E\):
\[
 \begin{cases}
 t/2-13/48+\lambda_1/3+2\lambda_2/3,&e\le1/8,\\
 t/2-9/32+5\lambda_1/6+\lambda_2/6,&1/8\le e\le1/4,\\
 t/2-11/32+\lambda_2,&e\ge1/4.
 \end{cases}
\]
In the first two lines use the equality sixth and twenty-fourth
moments respectively and interpolate the frequency second and
fourth moments. The last line uses the equality supremum,
the longer second and shorter fourth moments.
At \(t=9/20+3b/5\), the first branch of \(\lambda_j\)
dominates for both \(j\), since the lost sieve growth across
the length shortening is at least \(b/10\).
The first excess is zero, and the other two are nonpositive.
The last branch at \(t=2/3,\ b=41/144\) is zero.
The fixed \(S\)-part and sixth-power multipliers retain
their exact raw values on each fibre.

If the frequency sixth-power base is at least \(D^a\),
the product of \(j\) such bases divides the grouped base,
and summing its inverse-cube weights gives \(D^{-2ja}\).
The replacement \(\lambda_j\) is
\(-2a+\Delta_6(j(t-6a)_+)/(2j)\).
If the \(S\)-part has norm at least \(D^\sigma\), its product
has norm at least \(D^{j\sigma}\).
Summing the inverse square roots over this fixed prime support
costs \(D^{-j\sigma/2+\epsilon}\); use the Euler product
with any fixed positive exponent below \(1/2\).
Now \(\lambda_j=-\sigma/2+\Delta_6(j(t-\sigma)_+)/(2j)\).
At \(t=3/4\) these comparisons meet the target with
\(a=1/16\) and \(\sigma=3/10\).
Monotonicity gives smaller \(t\) and the original radial
decay sums the upper tail. All comparisons keep the
corresponding middle subtraction.

For the equality-wide sector, fix \(g,h,e_0,e_1,\ell/\ell_{\rm g}\).
The last factor lies on finitely many conjugates of primes
in \(S\), and its phase is retained.
The live dyad gives \(E_1Nh<14\).
On the remaining conjugate-good part the actual phase is
\(\nu(\ell_{\rm g})\overline{\nu(\bar\ell_{\rm g})}
 (Nu/\ell_{\rm g})_6\), including zeros at either orientation.
On \(T\le N\ell_{\rm g}<2T\) the rational conductor is
\(N\ell_{\rm g}\) and the squared coefficient mass is
of order at most \(1/T\) up to subpowers. Fixing the
second pivot, the oriented rational Gram bound has sieve
factor \(A+T^{5/3}\) in the integer variable \(Nv_1\);
equal-norm row multiplicity has divisor cost.
The sum over \(g\) is harmonic. Including the alternative
direct count gives normalized equality energy at most
\[
 \frac{D}{E^2}\left(1+\frac{D}{ER}\right)^{-8}
      \min\left(1,\frac{A+T^{5/3}}{AT}\right)
       (DH)^\epsilon\mathcal M_{24}(W)^2.
\]
For the full partner complement the four upper powers
of normalized energy are
\(3/4-e,\ 1-2e,\ 23/24-4e/3,\ 13/12-11e/6\).
For its wide coefficient the same bound follows by the
primal unit reduction, the exact small-prime compensation
and the physical sieve at \(D/E\); for the step use the
difference of the full and wide coefficients.
Cauchy requires the extra equality gain
\(\max(0,P(e)-1/2-8(1/4-e)_+)\), where \(P(e)\) is the
maximum of the four displayed powers.
This gives \(\theta(e)\) in Theorem~\ref{pca:joint-eqwide}.
The actual support \(ET\ll D\) ensures that the
\(T^{2/3}/A\) term does not obstruct this threshold.
The harmonic divisor sum finishes the full-complement bound.

For the distinct comparable selected rectangles, retain
the extra equality factor \((1+D/(ER))^{-7/2}\) after the
original two half radial powers. The coupled cutoff estimate
on the same insertion rectangles costs \(D^{1/17}\)
relative to \(R/E\). This estimate accepts the original
coefficient or its step and wide parts. At fixed ordered
\(m,n\) and, for a subtraction, fixed \(p\), the additional
wide tests are single-factor lower endpoints on the two
ordered arrays. The step tests are single-factor upper
endpoints. The exact common-prime projectors and actual
insertion dual retain these endpoints. Since
\((7/2)(1/4-111/476)=1/17\), the equality radial factor pays
that cost on the stated prefix. The selected middle subtraction
has its separate \(R/E\) bound. This proves the second
assertion of Theorem~\ref{pca:joint-eqwide}.
